\section{人：语言是为了泛化而发明的工具}

本章先看“人”的部分。姚顺雨的一个关键表达是：语言是人为了实现泛化而发明出来的工具，这一点比其他东西更本质。这个观点解释了他为什么从语言出发做 Agent：语言不仅是交流媒介，也是人类压缩经验、传递规则、描述目标、组织计划和跨场景迁移的工具。

\begin{figure}[H]
\centering
\includegraphics[width=0.94\textwidth]{language-generalization.png}
\caption{语言是泛化工具：人类用语言压缩经验，Agent 用语言连接任务和世界。自制概念图，依据 00:02:03--00:07:50 对谈内容整理。}
\end{figure}

\begin{knowledgebox}{读图：语言为什么适合做 Agent 的中间层}
语言把任务、规则、约束、工具说明和反馈都变成可组合的符号。模型通过语言可以理解人类目标，也可以调用代码、API、文档和其他 Agent。因此语言不是 Agent 的全部，但它是很强的通用接口。
\end{knowledgebox}

\subsection{非共识：早早押注语言 Agent}

上一节解释语言为什么是泛化工具，本节看这个判断如何变成研究下注。姚顺雨说自己一直有一个非共识：想要做 Agent，而且第一件事就是基于语言模型做 Agent。在当时，主流路线并不一定认为语言模型是通向 Agent 的最佳路径；但后来 GPT 系列显示，语言模型拥有强泛化和工具接口潜力，这个非共识变成了前沿主线。

\begin{importantbox}{Different Bet 的意义}
如果只复制已有霸主的主线，很难超越霸主。新的机会往往来自 different bet：不同任务、不同环境、不同 interface、不同产品形态。语言 Agent 在早期就是这样一种下注。
\end{importantbox}

\begin{knowledgebox}{为什么早期语言 Agent 是非共识}
早期语言模型看起来更像文本系统，而不是能行动的智能体。要把它看成 Agent，需要同时相信三件事：语言能表达任务，模型能跨任务泛化，外部工具可以把语言计划变成动作。这三件事在当时都不是显然共识。
\end{knowledgebox}

\begin{knowledgebox}{姚顺雨研究脉络的教学读法}
\begin{tabularx}{\textwidth}{p{0.20\textwidth}p{0.34\textwidth}X}
\toprule
阶段 & 关注问题 & 与 Agent 的关系 \\
\midrule
博士早期 & 用语言模型做简单游戏和任务 & 证明语言可以承载决策和状态。 \\
任务环境 & 寻找更真实、更有价值的 benchmark & 让 Agent 不只是刷小题。 \\
工具/代码 & InterCode、代码环境、数字世界 affordance & 让 Agent 获得可执行的手。 \\
下半场 & reward、multi-agent、interface、系统演化 & 从单模型能力进入系统能力。 \\
\bottomrule
\end{tabularx}
\end{knowledgebox}

\subsection{本章小结}

“人”的章节说明，Agent 研究不是凭空从模型能力里长出来的，而是来自对语言、人类泛化和任务表达的理解。语言之所以重要，是因为它把人类目标和机器行动接在一起。

